\section{Ley areal del corte triádico: capacidad \(81\log 3\), operador de área, espectro y entrelazador colectivo de los canales \(90/120\)}
\label{hap:sec:area-entrelazamiento-barbero-rev11}
\index{Barbero--Immirzi!operador de área}
\index{entropía modular!Hamiltoniano modular}

La identificación HMT del funcional \(\Gamma_{6\to8}\) con la coordenada
de Barbero--Immirzi se realiza sobre el mismo borde triádico que soporta la
incidencia (90/120). El interior y su hoja de borde poseen igual cardinalidad,
\[
 \mathcal B=(\mathbb Z/9\mathbb Z)^3,
 \quad |\mathcal B|=729,
 \qquad
 \partial\mathcal B=(\mathbb Z/27\mathbb Z)^2,
 \quad |\partial\mathcal B|=729,
\]
mediante el reagrupamiento de seis trits en tres coordenadas nonádicas o dos
coordenadas de orden veintisiete.

\begin{theorem}[Ley areal del corte triádico]
\label{hap:thm:ley-areal-corte-triadico-rev11}
En el grafo cúbico \(N\times N\times N\) con capacidad unitaria, el corte
mínimo entre dos caras opuestas tiene capacidad \(N^2\). Para \(N=9\),
\[
 \mathcal A_{\min}=81,
 \qquad S_{\rm tri}=\mathcal A_{\min}\log3=81\log3.
\]
\end{theorem}
\begin{proof}
Las \(N^2\) rectas paralelas a la dirección normal son caminos disjuntos en
aristas, por lo que todo corte las intercepta. Una capa transversal contiene
exactamente \(N^2\) aristas y alcanza la cota. Cada enlace triádico aporta
\(\log3\) a la entropía de un estado máximamente mezclado.
\end{proof}

En el toro de borde \(27\times27\), sea \(A\) el bloque canónico
\(9\times9\). Su frontera se descompone de manera exacta en
\begin{equation}
 \partial A=\partial A_{90}\sqcup\partial A_{120},
 \qquad |\partial A_{90}|=|\partial A_{120}|=18.
 \label{hap:eq:corte-90-120-rev11}
\end{equation}
Los enlaces horizontales representan el canal (90), y los verticales, el
canal (120). Esta descomposición convierte la pareja de incidencias en una
descomposición operatoria del borde.

La entropía areal \(S_{\rm tri}=81\log3\) pertenece al estado máximamente
mezclado de los enlaces del corte. La entropía bicapa
\(S_{\rm bi}(y)\) de la \cref{hap:def:entropia-bicapa-rev6} pertenece a la
distribución normalizada entre las dos hojas; ambos objetos conservan
dominios distintos.

Sea \(N=\operatorname{diag}(0,1,2)\) el operador número de un enlace trítico.
La realización areal declara la escala positiva
\[
 \theta_{\rm clk}=\frac{C^*}{6}\frac{\pi}{180},
 \qquad
 a_0=\frac{1+2\cos(10^\circ)}{3}\,\theta_{\rm clk},
 \qquad \gamma_{\rm HMT}=\Gamma_{6\to8}.
\]
La fórmula de \(a_0\) constituye la estructura de partida de esta
realización; fijada esa escala, el operador y su espectro se deducen
exactamente.

\begin{definition}[Operador de área HMT]
Sobre la hoja de borde
\(\mathcal H_{\partial}=\bigotimes_{e\in\partial A}\mathbb C^3\), se define
\begin{equation}
 \widehat{\mathcal A}_{\rm HMT}
 =\gamma_{\rm HMT}a_0\sum_{e\in\partial A}N_e.
 \label{hap:eq:operador-area-hmt-rev11}
\end{equation}
Su espectro en el corte canónico es
\[
 \operatorname{Spec}\widehat{\mathcal A}_{\rm HMT}
 =\{n\gamma_{\rm HMT}a_0:0\leq n\leq72\},
\]
con multiplicidades dadas por los coeficientes de
\((1+z+z^2)^{36}\). La prolongación por cortes compatibles produce la
familia protegida \(\mathcal A_n^{\rm prot}=n\gamma_{\rm HMT}a_0\).
\end{definition}

El entrelazador que fija la normalización se construye en el subespacio de
canales colectivos. Sean \(u_{90},u_{120}\) los vectores unitarios uniformes
de las incidencias hexádica y octádica, y sean \(v_{90},v_{120}\) los
vectores unitarios uniformes de los dos conjuntos de enlaces de
\eqref{hap:eq:corte-90-120-rev11}. El mapa
\begin{equation}
 \begin{aligned}
 \mathcal J_{6\to8}:{}
 &\operatorname{span}\{u_{90},u_{120}\}
 \longrightarrow
 \operatorname{span}\{v_{90},v_{120}\},\\
 &\mathcal J_{6\to8}u_{90}=v_{90},
 \qquad
 \mathcal J_{6\to8}u_{120}=v_{120}
 \end{aligned}
 \label{hap:eq:entrelazador-area-barbero-rev11}
\end{equation}
es una isometría entre estos subespacios colectivos. Si
\[
 D_{\rm inc}=90|u_{90}\rangle\langle u_{90}|
 +120|u_{120}\rangle\langle u_{120}|,
 \qquad
 D_{\rm area}=90|v_{90}\rangle\langle v_{90}|
 +120|v_{120}\rangle\langle v_{120}|,
\]
entonces
\[
 \boxed{\mathcal J_{6\to8}D_{\rm inc}
 =D_{\rm area}\mathcal J_{6\to8}.}
\]
En particular, la combinación primitiva \(12:-1\) actúa
con los mismos coeficientes antes y después de \(\mathcal J_{6\to8}\). Esta
igualdad fija el diccionario de normalización que inserta
\(\gamma_{\rm HMT}=\Gamma_{6\to8}\) en la conexión de
Ashtekar--Barbero.

